\section*{Bab \ref{ch:g10:energy-conservation} --- Energi: Bentuk dan Kekekalannya}

\begin{solution}{exo:g10:energy-conservation:1}
$E_k = \tfrac12 \times 1200 \times 25^2 = \qty{3.75e5}{J}$. Pada
\qty{36.1}{m/s}: $\tfrac12 \times 1200 \times 36.1^2 \approx
\qty{7.8e5}{J}$ --- laju yang $44\%$ lebih besar memberikan
$(130/90)^2 \approx 2.1$ kali \omterm{def:g10:energy-conservation:kinetic}{energi kinetiknya}: kuadratnya sedang
bekerja.
\end{solution}

\begin{solution}{exo:g10:energy-conservation:2}
$\Delta h = \qty{450}{m}$:
$E_p = 72 \times 9.81 \times 450 \approx \qty{3.2e5}{J}
= \num{3.2e5}/\num{3.6e6} \approx \qty{0.088}{kW.h}$, senilai sekitar
$0.022$ euro. Sehari mendaki setara dua sen listrik.
\end{solution}

\begin{solution}{exo:g10:energy-conservation:3}
$\qty{1}{kW.h} = \qty{3.6e6}{J}$;
$\qty{12}{W.h} = 12 \times \qty{3600}{J} \approx \qty{43}{kJ}$;
$\qty{45}{MJ} = 45/3.6 = \qty{12.5}{kW.h}$;
$\qty{10}{MJ} = 10/3.6 \approx \qty{2.8}{kW.h}$.
\end{solution}

\begin{solution}{exo:g10:energy-conservation:4}
$E = 2200 \times 150 = \qty{3.3e5}{J}
= \num{3.3e5}/\num{3.6e6} \approx \qty{0.092}{kW.h}$. Manusia:
$P = \num{1.0e7}/\num{86400} \approx \qty{116}{W}$.
\end{solution}

\begin{solution}{exo:g10:energy-conservation:5}
Dinamo: kimia (makanan pengendaranya) $\to$ kinetik (roda) $\to$
listrik (dinamo) $\to$ radiasi $+$ termal (lampu); kalornya bocor di
otot, sentuhan ban, dinamo, dan lampu. Pemanas bergas: kimia (gas)
$\to$ termal (air), bocor pada gas cerobongnya. Panel: radiasi
(Matahari) $\to$ listrik (panel) $\to$ kimia (baterai), kalornya bocor
di panel dan pengisi \omterm{def:g10:energy-conservation:power}{dayanya}.
\end{solution}

\begin{solution}{exo:g10:energy-conservation:6}
Yang dikonsumsi $2.0 \times 45 = \qty{90}{MJ}$; $\eta = 27/90 = 0.30$.
Kalornya: $90 - 27 = \qty{63}{MJ}$, yang dibawa pergi gas buang dan
radiatornya ke udara di sekitarnya.
\end{solution}

\begin{solution}{exo:g10:energy-conservation:7}
$E_k = E_p = 70 \times 9.81 \times 20 \approx \qty{1.37e4}{J}$;
$v = \sqrt{2 \times 13734/70} = \sqrt{392} \approx \qty{19.8}{m/s}
\approx \qty{71}{km/h}$.
\end{solution}

\begin{solution}{exo:g10:energy-conservation:8}
$41562 - 41236 = \qty{326}{kW.h}$; hariannya $326/30 \approx
\qty{10.9}{kW.h}$. Tagihannya: $326 \times 0.26 + 12 = 96.76$ euro.
\omterm{def:g10:energy-conservation:power}{Daya} rata-ratanya: $\qty{326}{kW.h}/\qty{720}{h} \approx \qty{453}{W}$.
\end{solution}

\begin{solution}{exo:g10:energy-conservation:9}
\emph{1.} Cahayanya: $0.05 \times 60 = \qty{3.0}{W}$; LED:
$\eta = 3/9 \approx 0.33$.
\emph{2.} Setahun ($\qty{1460}{h}$): bohlamnya
$60 \times 1460 = \qty{87.6}{kW.h}$, LED-nya $\qty{13.1}{kW.h}$;
penghematannya $\qty{74.5}{kW.h} \approx 19$ euro tiap bohlam tiap
tahun.
\end{solution}

\begin{solution}{exo:g10:energy-conservation:10}
\emph{1.} $P = \num{1.5e5} \times 9.81 \times 40 \approx
\qty{5.9e7}{W} = \qty{59}{MW}$.
\emph{2.} $0.90 \times 59 \approx \qty{53}{MW}$;
$\num{5.3e7}/525 \approx \num{1.0e5}$ rumah tangga.
\end{solution}

\begin{solution}{exo:g10:energy-conservation:11}
$v = \sqrt{2 \times 9.81 \times 0.050} \approx \qty{0.99}{m/s}$
(massanya saling menghapus: $\tfrac12 m v^2 = mgh$). \omterm{def:g10:energy-conservation:energy}{Energi} awalnya
$E_p = 0.20 \times 9.81 \times 0.050 \approx \qty{0.098}{J}$ berakhir
sebagai \omterm{def:g10:energy-conservation:forms}{energi termal} di udara dan pada porosnya: bandul itu tidak
terasing, dan jumlah seluruhnya --- ayunan ditambah kalornya ---
benar-benar kekal.
\end{solution}

\begin{solution}{exo:g10:energy-conservation:12}
\emph{1.} $\eta = 0.38 \times 0.94 \times 0.85 \approx 0.304$; yang
dikonsumsi $12/0.304 \approx \qty{39.5}{W.h}
= 39.5 \times 3600 \approx \qty{142}{kJ}$.
\emph{2.} $\num{1.42e5}/\num{2.5e7} \approx \qty{5.7}{g}$ batu bara
tiap pengisian; $\times 365 \approx \qty{2.1}{kg}$ tiap tahun.
\end{solution}

\begin{solution}{exo:g10:energy-conservation:13}
\emph{1.} $\num{5e9}/\num{3.6e6} \approx \qty{1400}{kW.h}$;
$1400/12.6 \approx 110$ hari.
\emph{2.} $P = \num{5e9}/\num{1e-4} = \qty{5e13}{W} = \qty{50}{TW}$
--- tujuh belas kali \omterm{def:g10:energy-conservation:power}{daya} listrik rata-rata umat manusia, selama
\qty{100}{\micro s}, di tempat yang tidak dapat diramalkan. Petir
adalah \omterm{def:g10:energy-conservation:power}{daya} yang besar tetapi \omterm{def:g10:energy-conservation:energy}{energi} yang kecil, yang diserahkan
terlalu cepat untuk berguna: tidak ada \omterm{def:g10:energy-conservation:transfer}{pengubah} yang sanggup menerima
lonjakan semacam itu.
\end{solution}

\begin{solution}{exo:g10:energy-conservation:14}
\emph{1.} $E_k = \tfrac12 \times 1300 \times 30.6^2 \approx
\qty{6.1e5}{J}$, yang diubah menjadi \omterm{def:g10:energy-conservation:forms}{energi termal} pada cakram dan
kampas remnya (lalu ke udara).
\emph{2.} $0.60 \times \qty{6.1e5}{J} \approx \qty{3.7e5}{J}$ tiap
perhentian; $\num{3.6e6}/\num{3.7e5} \approx 10$ perhentian tiap
\omterm{def:g10:energy-conservation:kwh}{kilowatt-jam}.
\end{solution}

\begin{solution}{exo:g10:energy-conservation:15}
$E_p = 85 \times 9.81 \times 1200 \approx \qty{1.0}{MJ}$; makanannya:
$1.0/0.25 = \qty{4.0}{MJ}$, yaitu $4.0/1.5 \approx 2.7$ mangkuk sereal
--- atau $4.0/45 \approx \qty{0.09}{kg}$ bensin. Satu pagi penuh usaha
yang jujur muat di dalam secangkir kopi bahan bakar: \omterm{def:g10:energy-conservation:forms}{energi kimia}
terpusat secara luar biasa.
\end{solution}

\begin{solution}{pb:g10:energy-conservation:1}
\textbf{1.} Lemari es $0.1 \times 24 = \qty{2.4}{kW.h}$; pemanas air
$2.4 \times 2.5 = \qty{6.0}{kW.h}$; oven \qty{2.0}{kW.h}; mesin cuci
\qty{1.0}{kW.h}; penerangan \qty{0.3}{kW.h}; layar \qty{0.9}{kW.h}.
\textbf{2.} Seluruhnya \qty{12.6}{kW.h} $= 12.6 \times 3.6 \approx
\qty{45}{MJ}$: satu \omterm{def:g10:orders-of-magnitude:unit}{kilogram} bensin, sampai ke \omterm{def:g10:energy-conservation:energy}{joulenya}.
\textbf{3.} $\qty{12600}{W.h}/\qty{24}{h} = \qty{525}{W}$ --- rumah itu
menyala pada seperempat sebuah ketel.
\textbf{4.} Pemanas airnya (\qty{6.0}{kW.h}, hampir separuhnya). Pada
\qty{1.25}{h}: pemanasnya \qty{3.0}{kW.h}, seluruhnya \qty{9.6}{kW.h}.
\textbf{5.} $15 \times 24 = \qty{0.36}{kW.h}$ tiap hari, sekitar
$2.8\%$ dari jumlah yang terukur meteran; setahun
$0.36 \times 365 \approx
\qty{131}{kW.h} \approx 33$ euro --- untuk
lampu yang tidak ditonton siapa pun.
\textbf{6.} $56604 - 56214 = \qty{390}{kW.h}$;
$390/30 = \qty{13.0}{kW.h}$ tiap hari: yaitu \qty{12.6}{kW.h} pada
auditnya ditambah \qty{0.36}{kW.h} lampu siaganya --- cocok.
\textbf{7.} $390 \times 0.25 + 14 = 111.50$ euro.
\textbf{8.} Jaringan listrik: $0.25/3.6 \approx 0.07$ euro tiap
megajoule. Bensin: $1.80/34 \approx 0.05$ euro tiap megajoule.
\textbf{9.} Mangkuknya: $0.50/1.5 \approx 0.33$ euro tiap megajoule ---
lima kali jaringan listriknya. Dari yang termurah: bensin, jaringan
listrik, lalu makanan, jauh tertinggal.
\textbf{10.} $0.20 \times 1000 = \qty{200}{W}$ tiap \omterm{def:g10:orders-of-magnitude:unit}{meter} persegi.
\textbf{11.} $200 \times 4.0 = \qty{800}{W.h} = \qty{0.80}{kW.h}$ tiap
\omterm{def:g10:orders-of-magnitude:unit}{meter} persegi tiap hari.
\textbf{12.} $12.6/0.80 = 15.75 \approx \qty{16}{m^2}$: muat pada atap
\qty{20}{m^2} itu.
\textbf{13.} $16 \times 0.2 \times 1.5 = \qty{4.8}{kW.h}$: sekitar
$38\%$ kebutuhannya; jaringan listriknya (atau baterai yang terisi pada
bulan yang lebih baik) harus memasok sisanya.
\textbf{14.} Satu \omterm{def:g10:orders-of-magnitude:unit}{meter} persegi menghasilkan $0.80 \times 365 \approx
\qty{292}{kW.h}$ tiap tahun; pengembaliannya $700/292 \approx 2.4$
tahun --- singkat dibandingkan umur panel yang 25 tahun.
\textbf{15.} $\qty{12.6}{kW.h}/\qty{0.1}{kW} = \qty{126}{h}$: lebih
dari lima orang ($126/24 = 5.25$) yang mengayuh siang dan malam.
\textbf{16.} Makanannya $= 12.6/0.25 = \qty{50.4}{kW.h} \approx
\qty{181}{MJ}$, yaitu $181/1.5 \approx 121$ mangkuk sereal.
\textbf{17.} $121 \times 0.50 \approx 60$ euro, berbanding
$12.6 \times 0.25 \approx 3.15$ euro dari jaringan listrik: tenaga
kayuh berharga sekitar $19$ kali lipat --- sebelum membayar para
pengayuhnya.
\textbf{18.} $5000/45 \approx 110$ hari rumah tangga. Tetapi
$\num{5e9}/\num{1e-4} = \qty{5e13}{W}$ selama \qty{100}{\micro s}, di
tempat yang acak: tidak ada \omterm{def:g10:energy-conservation:transfer}{pengubah} atau baterai yang sanggup minum
dari selang semacam itu. Sebuah rumah memerlukan aliran tetap
\qty{525}{W}, bukan sebuah lonjakan.
\textbf{19.} Bukan \omterm{def:g10:energy-conservation:energy}{joulenya} --- \omterm{def:g10:energy-conservation:energy}{joule} bersifat kekal dan serupa. Kamu
membayar \emph{bentuk dan ketersediaannya}: \omterm{def:g10:energy-conservation:energy}{energi} yang terpusat, dapat
disimpan dan tersedia sewaktu-waktu berharga sebesar seluruh \omterm{def:g10:energy-conservation:transfer}{rantai
pengubahan} yang membuatnya demikian.
\textbf{20.} Satu hari keluarga itu $= \qty{12.6}{kW.h} \approx
\qty{45}{MJ}$: satu \omterm{def:g10:orders-of-magnitude:unit}{kilogram} bensin, enam belas \omterm{def:g10:orders-of-magnitude:unit}{meter} persegi atap yang
tersinari matahari, atau lima pengayuh sepanjang siang dan malam.
Persediaannya adalah \omterm{def:g10:energy-conservation:energy}{energi} (\omterm{def:g10:energy-conservation:energy}{joule}, \omterm{def:g10:energy-conservation:kwh}{kilowatt-jam}); alirannya adalah
\omterm{def:g10:energy-conservation:power}{daya} (\omterm{def:g10:energy-conservation:power}{watt}) --- bak mandi dan kerannya, jangan lagi dicampuradukkan.
\end{solution}
